\documentclass[12pt]{article}
\usepackage[utf8]{inputenc}
\usepackage[T1]{fontenc}
\usepackage{hyperref}
\title{Miqie Language Resources for K-12 Curriculum Development}
\author{Compiled by Claude Sonnet 5.5 (claude-sonnet-5-5) for Andr\'as Kornai}
\date{}
\begin{document}
\maketitle

\section{Introduction}
Miqie (also called Micha; ISO 639-3 \texttt{yiq}, Glottocode miqi1235) is a Loloish (Ngwi) language of north-central Yunnan, China. It is spoken in Wuding, Luquan and Fumin counties and belongs to the Sino-Tibetan family. Speaker estimates range widely, from about 8,000 to 30,000. Sources report no established orthography. Do not confuse it with Lalo, which has a similar-sounding autonym.

\section{Primary Educational Resources}
No primers, textbooks or school programmes were found.

\section{Linguistic and Reference Materials}
\begin{itemize}
\item Micha language, Wikipedia: \url{https://en.wikipedia.org/wiki/Micha_language}
\item Endangered Languages Catalogue entry: \url{https://endangeredlanguages.com/lang/3191?hl=zh-cn}
\item Endangered Languages Catalogue bibliography (it cites Gao 2015 in \emph{Language Documentation \&amp; Conservation}): \url{https://endangeredlanguages.com/lang/3191/bibliography?hl=es}
\item Verbix Endangered Languages entry: \url{https://docs.verbix.com/EndangeredLanguages/MiqieYi}
\item Chinese Wikipedia article: \url{https://zh.wikipedia.com/wiki/%E5%AF%86%E5%AF%9F%E8%AF%AD}
\end{itemize}

\section{Religious and Community Resources}
None found.

\section{Digital and Audio Resources}
No audio, video or corpus resources were found in the searches.

\section{Curriculum Development Context}
Documentation is thin and no writing system is reported. A curriculum would first need an orthography and community materials. Gao (2015) is the main scholarly lead and should be located and read in full.
\end{document}